% 原文 pp.240–245（印刷 pp.1358–1363）
% 第 17 讲第 2 部分：本文件承接第 1 部分末尾的 §17.2.4，故不写 \section，直接从正文起。
% 说明：p.240–245 上原文的展示公式均未印编号（右边界无编号栏），故一律用 \[...\]；
% p.245 参考文献提到的「拉格朗日量 (17.1)」当指第 17 讲靠前处（第 1 部分）为 $\mathcal{N}=2$ $U(1)$
% 矢量多重态拉格朗日量所加的编号，其形式即本文件 p.241 上重录（未编号）的那一式。

% ===== 原文 p.240 =====
% 承接 §17.2.4（椭圆纤维化的全纯二形式 $\eta$）：本页从上一页句子中间续起。

……的 $a^\alpha$。尽管如此，二形式 $da^\alpha\wedge dz$ 确实能够下降（descend）为 $\mathcal{E}^\alpha$ 上一个定义良好的二形式。显然，$\mathcal{E}|_{V^\alpha}$ 上满足上面第一个积分方程的全纯二形式，只能是

\[
  \eta^\alpha = -da^\alpha\wedge dz .
\]

% [Tutor Note: 这里的「显然」实际用了两个事实。(i) 在 $\mathcal{E}|_{V^\alpha}$ 上，全纯二形式必呈 $f(a)\,da^\alpha\wedge dz$ 的形状：可用的全纯微分只有底空间拉回的 $da^\alpha$ 与沿纤维的 $dz$，而 $da\wedge da=dz\wedge dz=0$，故这类形式在「全纯函数倍」的意义下一维（即 $H^0(E,\Omega^1)=\mathbf{C}\,dz$ 的一维性）。(ii) 上一页的「第一个积分方程」（等价地 p.242 的 $\int_{E_1}\lambda=a^\alpha$，两边取 $d$ 得 $\int_{E_1^\alpha}\eta^\alpha=da^\alpha$）把这个函数钉住：在 $E=\mathbf{C}/L(1,\tau)$ 上 $\int_{E_1}dz=1$，代入 $f(a)\,da\wedge dz$ 即得 $f=-1$，负号来自 $\eta^\alpha=-da^\alpha\wedge dz$ 的定向约定。所以「唯一性」是归一化条件的推论，而非同调层面的显然性。另注意 $dz$ 并不是 $\mathcal{E}$ 上的整体形式（$z$ 只在每个 $V^\alpha$ 上是局部坐标），能被粘成整体形式的只是 $\eta$ 本身。]

这个二形式 $\eta^\alpha$ 显然也满足第二个积分方程。现在我们需要说明：这些局部二形式能够拼接起来，给出一个整体形式 $\eta$。但这由局部形式的唯一性立刻可见。

现在我们来计算 $\pi_*(\eta\wedge\bar\eta)$，其中 $\pi$ 是投影 $\mathcal{E}\to U$。让我们把注意力限制在 $\mathcal{E}|_{V^\alpha}$ 上。由我们的局部描述有

\[
  \pi_*(\eta\wedge\bar\eta)=\int_{\text{fibers}}d\lambda^\alpha\wedge d\bar\lambda^\alpha .
\]

由 $d\lambda$ 在 $E_1^\alpha$ 与 $E_2^\alpha$ 上的积分公式，以及 $E_1^\alpha\cdot E_2^\alpha=1$（在 $E(u)$ 的同调中成立）这一事实，立即得到结果为

\[
  \pi_*(\eta\wedge\bar\eta)
  =da^\alpha\wedge d\bar a^\alpha\,\bar\tau^\alpha
   +da^\alpha\wedge d\bar a^\alpha\,\tau^\alpha
  =(\tau^\alpha-\bar\tau^\alpha)\,da^\alpha\wedge d\bar a^\alpha
  =8\pi\,\Omega|_{V^\alpha}.
\]

% [Tutor Note: 中间这一步的两个被加项在印刷本上都写作 $da^\alpha\wedge d\bar a^\alpha$，但若二者真按同一次序，其和应为 $(\tau^\alpha+\bar\tau^\alpha)da^\alpha\wedge d\bar a^\alpha$，与紧接着写出的 $(\tau^\alpha-\bar\tau^\alpha)$ 不符。要得到印刷结果，第二项必须理解为次序相反的 $d\bar a^\alpha\wedge da^\alpha\,\tau^\alpha=-da^\alpha\wedge d\bar a^\alpha\,\tau^\alpha$——即 $1$-形式的交换带来一个负号，而系数 $\tau^\alpha$ 与 $da^\alpha$ 都是全纯的、彼此可交换，不会自行产生符号。抄录或排印时最容易在此丢一个负号，读者务必按 $\tau-\bar\tau=2i\,\mathrm{Im}\,\tau$ 自行核对。另：$8\pi$ 这个因子来自 $\lambda$ 沿实二维纤维的积分（$\int_E dz\wedge d\bar z=2i\,\mathrm{Im}\,\tau$ 型的面积因子）与规范动力学项 $\frac{1}{4\pi}\mathrm{Im}\,\tau\,F^2$ 的归一化之积，故 $\Omega$ 与 $\mathrm{Im}\,\tau\,da\wedge d\bar a$ 之间相差一个常数因子，不可把 $\Omega$ 直接当成 $\mathrm{Im}\,\tau\,da\wedge d\bar a$。]

$\Omega$ 一如前述，是 $U$ 上度量的 Kähler 形式。

由于二形式 $\eta$ 是全纯的，它在任何纤维上、以及在椭圆纤维化的截面 $\sigma$ 上的积分都平凡地为 $0$。因此，它的同调类由 $\eta$ 沿「位于 $U$ 中圆周上方的二维环面」的积分所决定。

现在我们要说明：$\eta$ 的 de Rham 上同调类，决定了在用各种 $(a_D,a,\tau)$ 数组给出的局部描述之间转换时所需的自同构究竟落在 $SL(2,\mathbf{Z})$ 中，还是落在扩充群 $ISL(2,\mathbf{Z})$ 中。我们断言：若 $\eta$ 的 de Rham 上同调类平凡，则我们用到的自同构落在 $SL(2,\mathbf{Z})$ 中。要点如下。我们已经知道 $da_D,da$ 的变化由 $SL(2,\mathbf{Z})$ 中的元素给出；我们还需要计算 $a_D,a$ 的变化。若我们沿 $U$ 中任一回路 $\gamma$ 绕行，且 $E_1$ 在该回路绕行后回到自身，则各循环 $E_1(\gamma(t))$ 拼成一个环面，而 $\eta$ 沿该环面的积分等于 $\int_\gamma da$。若此积分为零，则 $a$ 可以沿此回路解析延拓为单值函数。类似的论证适用于 $a_D$。于是，若 $\eta$ 的所有周期都为零，则 $a,a_D$ 按 $SL(2,\mathbf{Z})$ 变换。

一般而言，像 $\eta$ 这样的全纯二形式，其 de Rham 上同调类平凡，当且仅当 $\mathcal{E}$ 上存在整体的第二类微分（differential of the second kind）$\lambda$ 使 $d\lambda=\eta$。在更一般的情形下，我们可以把全纯二形式 $\eta$ 写成 $d\lambda$，其中 $\lambda$ 是 $\mathcal{E}$ 上某个整体的第三类微分（即 $\lambda$ 是 $\mathcal{E}$ 上的亚纯一形式，沿其极点轨迹的留点为常数）。结果表明，在我们对纯 $SU(2)$ 规范理论的应用中，各种描述之间的自同构都落在 $SL(2,\mathbf{Z})$ 中，因此我们无需考虑更一般的情形。这一事实的证明是下一小节的主题。

\subsection{BPS 形式}
% 原文 §17.2.5：The BPS formalism

我们这里的目标是说明：若我们一直考虑的低能有效 $U(1)$ 规范理论是纯 $\mathcal{N}=2$ 超对称 $SU(2)$ 规范理论的低能极限，则 $U(1)$ 理论各种以 $(a_D,a,\tau)$ 表述的局部描述之间的变换落在 $SL(2,\mathbf{Z})$ 中。我们将借助 BPS 形式（BPS formalism）来确立这一点。在任何 $U(1)$ 理论中，我们都有两个守恒电荷，即

% ===== 原文 p.241 =====

电荷与磁荷：

\[
  Q_e=\frac{1}{4\pi}\int_{S^2}*F,\qquad
  Q_m=\frac{1}{4\pi}\int_{S^2}F .
\]

这里，我们是在闵可夫斯基四维空间中某个时刻的无穷远球面上作积分。在四维空间上的纯 $U(1)$ 理论中，这两个荷都是零。但若我们的 $U(1)$ 理论是某个 $SU(2)$ 理论的低能极限，则这些荷可以非零。回忆第 8 讲 §5：我们的 $SU(2)$ 理论有一个 $U(1)$ 对称性，其守恒荷——至少在 $\tau$ 为常数时——为

\[
  N_e=Q_e-\frac{\theta Q_m}{2\pi}.
\]

由于它生成一个 $U(1)$ 对称性，这个算符的本征值为整数。还有另一个守恒荷，即磁荷，由

\[
  N_m=Q_m
\]

给出。它由无穷远处线丛的第一 Chern 类决定，因此显然也具有整本征值。既然两者的本征值都是整数，这两个荷就都联系着 $U(1)$ 对称性。当然，正如我们在第 8 讲 §5 中看到的，这两个荷在对偶下互相交换。

% [Tutor Note: $N_e=Q_e-\theta Q_m/2\pi$ 是 Witten 效应（Witten effect）的表达式：$\theta$ 角使磁单极子带上分数电荷，因而「整数荷」并非 $Q_e$ 而是 $Q_e-\frac{\theta}{2\pi}Q_m$。这里 $Q_m=\frac{1}{4\pi}\int_{S^2}F$ 的归一化使 't Hooft--Polyakov 磁单极子的 $Q_m$ 为整数（Dirac 量子化已取 $e=1$ 的单位），于是 $\theta Q_m/2\pi$ 随 $\theta\to\theta+2\pi$ 平移整数，$N_e$ 的整性得以保持——这正是 $\theta$ 的 $2\pi$ 周期性在荷谱上的体现。切勿把 $N_e$ 与 $Q_e$ 混用：在 $\theta\ne0$ 时二者相差一个分数，而由指标定理保证整性的是 $N_e$。]

$\mathcal{N}=2$ 超对称生成元 $Q^i_\alpha$（$i=1,2$）满足超 Poincaré 群的（反）关系，其中包括

\[
  \{Q^i_\alpha,Q^j_{\dot\alpha}\}=\delta^i{}_j\,P_{\alpha\dot\alpha},
\]

其中 $P_{\alpha\dot\alpha}$ 是沿方向 $\{\alpha\dot\alpha\}$ 的平移生成元。此外还有

\[
  \{Q^i_\alpha,Q^j_\beta\}=\epsilon^{ij}\epsilon_{\alpha\beta}\,S,
\]

其中 $S$ 是理论的中心荷（central charge）；在每一个超多重态中，$S$ 就是乘以一个标量。

% [Tutor Note: 这两条反交换关系里指标的位置是初读时最易出错之处：$\{Q^i_\alpha,Q^j_{\dot\alpha}\}=\delta^i{}_jP_{\alpha\dot\alpha}$ 中 $i$ 在上、$j$ 在下（一个上标一个下标才可缩并），而 $\{Q^i_\alpha,Q^j_\beta\}=\epsilon^{ij}\epsilon_{\alpha\beta}S$ 是 $\mathcal{N}=2$ 的中心扩张项，两个 $SO(2)_R$ 指标都取上标、由 $\epsilon^{ij}$ 缩并。许多文献把中心荷写成 $Z$ 并取相反的符号（$\{Q,Q\}=\epsilon\epsilon\,\bar Z$），此时 BPS 界仍是 $M\ge|Z|$，但 $Z$ 与本文 $S$ 差一个复共轭；又因 $\epsilon^{12}=1$、$\epsilon_{12}=-1$ 的约定不同，中心荷项的整体符号还会再翻一次。凡跨文献比较「BPS 态」时，务必先核对这一符号约定，否则会把 $aN_e+a_DN_m$ 误写成其复共轭。]

一般而言，任何态的质量都受 BPS 界（BPS bound）的限制，即

\[
  M\ge|S| .
\]

对 BPS 饱和态——即使上述不等式取等号的态——人们往往能获得很好的理解。

在我们的情形中，我们希望计算中心荷 $S$。为此回忆：拉格朗日量（$\tau$ 为常数时）为

\[
  \mathcal{L}=\frac{1}{4\pi}\int d^4x\,d^2\theta\;i\tau W_\alpha W^\alpha
  +\text{c.c.}
  +\int d^4x\,d^4\theta\;\mathrm{Im}\,\tau\,\bar\Phi\Phi .
\]

由此可知（见第 8 讲 §5）：对该 $\tau$ 值有

\[
  S=aN_e+a_DN_m,
\]

其中 $\Phi=a+\cdots$。按量纲论证，即使 $\tau$ 不是常数，此式也成立；含 $a$ 与 $a_D$ 的导数项在量纲上不可能出现。其解释是：低能理论的不同拉格朗日描述使用不同的 $a_D$ 与 $a$ 值；在每一种描述中，$S$ 都按上式依赖于 $a_D$ 与 $a$ 的取值。回忆第 8 讲 §5：当我们取理论的不同表示时，$(N_m,N_e)$ 会按

% ===== 原文 p.242 =====

$SL(2,\mathbf{Z})$ 变换。由于 $S$ 必须保持不变，我们看出：若

\[
  \begin{pmatrix} a_D\\ a\end{pmatrix}
  \longmapsto M\cdot\begin{pmatrix} a_D\\ a\end{pmatrix},
\]

则

\[
  \bigl(N_m\quad N_e\bigr)\longmapsto
  \bigl(N_m\quad N_e\bigr)\cdot M^{-1}
\]

对任何 $M\in SL(2,\mathbf{Z})$ 成立。

% [Tutor Note: 这两式的行列排布必须一起记：$S=aN_e+a_DN_m$ 可统一写成矩阵乘积 $S=(N_m\ N_e)\binom{a_D}{a}$，所以 $a$ 的列向量左乘 $M$ 时，荷的行向量必须右乘 $M^{-1}$（而不是 $M^{-\mathsf{T}}$ 或 $M^{\mathsf{T}}$），乘积才不变。若把荷写成列向量 $(N_e,N_m)^{\mathsf T}$，则变换矩阵要整体转置；若误把行向量写成 $(N_e\ N_m)$ 而仍用 $M^{-1}$，就会把 $N_m$ 和 $N_e$ 的角色互换，直接毁掉后面 $N_e$ 的整性论证。]

首先要指出的重要一点是：我们无法通过改变 $(N_m,N_e)$ 来补偿「给 $\binom{a_D}{a}$ 加上一个常数」这一操作，因为在这个理论中没有别的东西（没有别的守恒荷）可供 $N_m$ 与 $N_e$ 变换过去。这就说明了我们在上一小节末尾所断言的事，即：在只含 $\mathcal{N}=2$ 矢量多重态的 $SU(2)$ 规范理论中，$U(1)$ 理论各种以 $(a_D,a,\tau)$ 表述的表示之间的自同构落在 $SL(2,\mathbf{Z})$ 中。因此，在此情形下我们可以在 $\mathcal{E}$ 上构造一个第二类微分 $\lambda$，使得对每个 $\alpha$ 都有

\[
  \int_{E_1}\lambda=a^\alpha,\qquad
  \int_{E_2}\lambda=a_D^\alpha
\]

在 $\mathcal{E}|_{V^\alpha}$ 上成立。

更一般地，若 $SU(2)$ 理论具有额外的对称性（例如当我们加入超多重态时可能出现这种情况），则中心荷的形式可能变成

\[
  S=aN_e+a_DN_m+W,
\]

其中 $W$ 是某个别的守恒荷。在此情形下，就有可能补偿 $(a_D,a)$ 中的仿射平移（affine translation）。也就是说，我们可以有

\[
  \begin{pmatrix} a_D\\ a\\ 1\end{pmatrix}
  \longmapsto
  \begin{pmatrix}
    M & \binom{c_D}{c}\\[4pt]
    (0\quad 0) & 1
  \end{pmatrix}
  \cdot
  \begin{pmatrix} a_D\\ a\\ 1\end{pmatrix},
\]

以及

\[
  \bigl(N_m\quad N_e\quad W\bigr)\longmapsto
  \bigl(N_m\quad N_e\quad W\bigr)\cdot
  \begin{pmatrix}
    M^{-1} & \binom{*}{*}\\[4pt]
    (0\quad 0) & 1
  \end{pmatrix}.
\]

在所有这些变换中都不出现 $\tau$ 的导数。事实上，我们所做的是先对 $\tau$ 为常数算出变换律，然后以显然的方式把它们推广到 $\tau$ 全纯的情形。我们之所以知道这一机制可行，是因为仅凭量纲论证就可以断定：公式中不可能出现 $\tau$ 的导数。

% [Tutor Note: 这里的「仿射平移」矩阵属于 $\mathbb{Z}$ 上的 $ISL(2,\mathbf{Z})$，但作者没有写出整性条件，也没有写出右上角那一列究竟是什么。要使 $S=(N_m\ N_e\ W)\cdot(a_D,a,1)^{\mathsf T}$ 在 $(a_D,a)\mapsto M(a_D,a)+\binom{c_D}{c}$ 下不变，必须同时令 $(N_m\ N_e)\mapsto(N_m\ N_e)M^{-1}$ 且 $W'=W-(N_m\ N_e)M^{-1}\binom{c_D}{c}$；换句话说，印刷本中写成 $\binom{*}{*}$ 的那一列由不变性唯一确定为 $-M^{-1}\binom{c_D}{c}$（一个二分量列向量），并不是可以任取的。另一个易错点：只有 $S$ 中真的多了 $W$ 这样一项（即存在能与常数平移配对的新荷）时，才能吸收常数平移；纯 $\mathcal{N}=2$ 矢量多重态之所以没有这种自由度，恰恰因为荷的个数不够——这正是前面「没有别的东西可供变换」一句的矩阵表述。]

我们将主要限于没有额外对称性、且低能理论各种表示之间的变换落在 $SL(2,\mathbf{Z})$ 中的情形。但在这里以及第 19 讲中，我们会就带有额外超多重态、确实会遇到仿射平移的情形作若干评注。

关于 $SU(2)$ 理论中心荷的形式，我们已经给出了一种推导。这里再给出同一公式的另一种解释（至少是其一部分）。让我们考虑带一个带电 $\mathcal{N}=2$ 超多重态的 $U(1)$ 规范理论。在 $\mathcal{N}=1$ 的语言中，这个超多重态是一对手征超场 $(T,\tilde T)$。

% ===== 原文 p.243 =====

超对称要求这两个荷必须相反；我们将其记为 $\pm n_e$。$\mathcal{N}=2$ 超对称还要求：在 $\mathcal{N}=1$ 的描述中，除可能的裸质量外，超势中依赖于 $T,\tilde T$ 的部分恰为

\[
  \tilde W=n_e\Phi T\tilde T .
\]

由于 $\Phi=a+\cdots$，我们看到 $(T,\tilde T)$ 这一对的质量为 $an_e$。这个有质量的超多重态是 BPS 的（其态的重数太小，不可能是别的），因此该多重态的质量为 $an_e$ 这一事实意味着中心荷必须含有 $aN_e$ 这一项。（$T$ 场是 $N_e$ 的本征态，本征值为 $n_e$；它也是 $N_m$ 的本征态，本征值为 $0$ —— 这正是我们无法用这种办法探测中心荷中 $a_DN_m$ 项的原因。）

% [Tutor Note: 作者的推理链是：$\mathcal{N}=1$ 超势中的质量项 $n_e\Phi T\tilde T$ 在 $\Phi$ 取其标量分量 $a$ 处给出质量 $m_T=n_ea$；另一方面 BPS 公式给出 $M=|S|=|n_ea+n_ma_D|$，而 $T$ 的 $N_m=0$，故 $S$ 必须含 $an_e$ 项。括号里「重数太小」是多重态计数：非 BPS 的长多重态所含态数是短（BPS）多重态的两倍，而一对带电手征超场恰好只提供短多重态那么多的态，凑不出长多重态所需的额外伙伴（超势中也没有这样的场），故它只能是短的。这个「数态」判据在 $\mathcal{N}=2$ 中是判定 BPS 的常用手段。]

若我们为这个多重态加上裸质量，超势的相关部分变成

\[
  \tilde W=(n_e\Phi+m)T\tilde T,
\]

其中 $m$ 为常数。此时质量为 $M=n_ea+m$，而多出的常数意味着中心荷实际上是 $S=an_e+a_DN_m+mW$，其中 $W$ 是一个新的守恒量，它在 $T$ 与 $\tilde T$ 上取 $1$ 或 $-1$（在矢量多重态上则平凡作用）。由于 $S$ 中出现了这个额外项，单值性（monodromy）最终会包含仿射平移。细节请参看文献。

% [Tutor Note: 这里的 $W$ 与本讲前面 $W_\alpha$（场强超场）、以及超势 $\tilde W$ 完全无关，只是字母重复使用；原文用 $\tilde W$ 与 $W$ 相区分，译文保留之，读者仍须按上下文辨别。物理上 $W$ 是 $T,\tilde T$ 的相对 $U(1)$ 荷算符（在 $T$ 上取 $+1$、$\tilde T$ 上取 $-1$）：裸质量项 $mT\tilde T$ 之所以在中心荷中留下 $mW$，是因为质量可看作该 $U(1)$ 的荷移动；而一旦 $S$ 中出现 $W$，$(a_D,a)$ 的常数平移便可由 $W$ 的重新定义吸收，$SL(2,\mathbf{Z})$ 因而扩充为 $ISL(2,\mathbf{Z})$——这正是「单值性最终包含仿射平移」的含义。]

\subsection{BPS 谱的跳变}
% 原文 §17.2.6：Jumping of the BPS spectrum

让我们回到由 $a,a_D$ 局部给出的低能 $\mathcal{N}=2$ 超对称 $U(1)$ 规范理论。在此类理论中，BPS 质量不等式为

\[
  M\ge|n_ea+n_ma_D| .
\]

由于 $n_e$ 与 $n_m$ 是整数，这个界是好的：其含义是，只要 $a$ 与 $a_D$ 在 $\mathbf{C}$ 中生成一个格（lattice），就存在一个正常数作为 $M$ 的下界（对 $n_e,n_m$ 非零的态而言）。

当 $a$ 与 $a_D$ 生成的格退化时——也就是当 $a$ 与 $a_D$ 在 $\mathbf{C}$ 中共线时——会发生什么？在这类点上，联系不同格点上各态质量的三角不等式不再是严格的。这就允许态衰变为一对更轻的态，而当 $a$ 与 $a_D$ 共线时，这对更轻的态与它简并。BPS 态就这样可以消失（或出现），这正是 BPS 谱跳变（jumping）的来源。

一般地，我们预期这发生在 $U$ 中的一条实曲线上。（注意：由 $(1,\tau)$ 生成的格并未退化，椭圆曲线本身没有任何毛病；退化的其实是由 $a$ 与 $a_D$ 生成的格。）这一现象与 Vafa 讲演中二维 $\mathcal{N}=2$ 超对称理论的情形相似。在那里，BPS 态是孤子（即我们在第 12.5 讲引入的那些孤子），而跳变发生在不同孤子的中心荷取值共线之时。在那些二维模型中，这发生在超势 $W$ 在复平面上有三个共线的临界点时。在我们今天讨论的四维例子中，我们将得到关于 $U$ 中 BPS 谱发生跳变的那条曲线的明确信息。

% ===== 原文 p.244 =====

\subsection{关于 $SU(2)$ 规范理论应用的进一步讨论}
% 原文 §17.2.7：More about the application to SU(2) gauge theory

正如我们已多次顺带提到的，我们将研究一个不含物质超多重态的 $\mathcal{N}=2$ 超对称 $SU(2)$ 规范理论。我们的做法是把该理论当作 $\mathcal{N}=1$ 超对称理论来研究，其中 $\Phi$ 是取伴随表示值的（adjoint-valued）手征超场。微观（即高能）拉格朗日量为：

\[
  \mathcal{L}_{\text{mic}}
  =\frac{1}{4\pi}\int d^4x\,d^2\theta\;\tau_0(\Lambda_0)\Tr W_\alpha W^\alpha
  +\text{c.c.}
  +\frac{1}{4\pi}\int d^4x\,d^4\theta\;\mathrm{Im}\,\tau_0(\Lambda_0)\Tr\bar\Phi\Phi,
\]

其中 $W$ 是闵可夫斯基空间上主 $SU(2)$ 丛的超联络（superconnection）$\mathcal{A}$ 之超曲率的一部分，而 $\Phi=\phi+\theta\psi+\cdots$，$\phi$ 是取值于伴随丛复化中的标量场。$\Lambda_0$ 是我们为正规化 $SU(2)$ 理论而选取的紫外动量截断。此外，$\tau_0$ 是（关于场的）常数，它依赖于 $\Lambda_0$ 的取值；当 $\Lambda_0\to\infty$ 时须调整 $\tau_0$，使理论有极限。

经典真空的空间由这样的位形给出：平凡联络，加上 $\mathbf{C}\otimes_{\mathbf{R}}\mathfrak{su}(2)$ 中使势为零的常（协变常）截面 $\phi$，再模去规范等价。此拉格朗日量的势是矩映射的模平方：

\[
  V=\frac{\mathrm{Im}\,\tau_0}{4\pi}\int d^4x\;\Tr\bigl([\phi,\bar\phi]^2\bigr).
\]

若 $\phi$ 是势函数的零点，则显然 $\phi$ 与 $\bar\phi$ 对易。写 $\phi=\gamma+i\delta$，其中 $\gamma,\delta\in\mathfrak{su}(2)$，我们看到「$\phi$ 与 $\bar\phi$ 对易」这一条件无非就是 $\gamma$ 与 $\delta$ 对易。这意味着 $\gamma$ 与 $\delta$ 可以同时对角化。也就是说，模去规范变换后有

\[
  \phi=\frac{1}{\sqrt2}
  \begin{pmatrix} a & 0\\ 0 & -a\end{pmatrix}
\]

其中 $a\in\mathbf{C}$。还有进一步的共轭，即 Weyl 群的作用；在此情形 Weyl 群为 $\mathbf{Z}/2\mathbf{Z}$，其作用为 $a\mapsto-a$。因此，经典真空的模空间是由复坐标

\[
  u=\Tr\phi^2=a^2
\]

参数化的复平面。

在经典理论中，当 $u\ne0$ 时，$SU(2)$ 规范对称性通过 Higgs 机制破缺为 $U(1)$，超场的许多分量获得质量。低能理论正是我们今天所讨论的那一类：其中的轻场构成一个具有 $\mathcal{N}=2$ 超对称的 $U(1)$ 规范多重态。而在例外点 $u=0$ 处，低能有效理论仍是不破缺的 $SU(2)$ 规范理论。

% [Tutor Note: $\phi=\frac{1}{\sqrt2}\mathrm{diag}(a,-a)$ 与 $u=\Tr\phi^2=a^2$ 这一对公式里，$\frac{1}{\sqrt2}$ 的因子常被抄错（写成 $\frac12$ 或干脆丢掉），后果是 $u$ 与 $a$ 的关系变成 $u=2a^2$ 或 $u=a^2/2$，进而使下一讲 Seiberg--Witten 解中奇点的位置（正比于 $\Lambda^2$）整体错一个因子。可按 $\Tr\phi^2=\frac12(a^2+(-a)^2)=a^2$ 自查。另请注意：$\phi$ 取值于 $\mathfrak{su}(2)$ 的复化，$\bar\phi$ 与 $\phi$ 不是独立变量；条件 $[\phi,\bar\phi]=0$ 是 $D$ 项（矩映射）为零，而非 $F$ 项条件。]

让我们结合今天关于低能有效 $U(1)$ 理论的讨论来考察它。正如刚才所解释的，经典真空空间是复平面——即 $u$ 平面。让我们考察 $u$ 平面上的度量。此平面上除原点外的所有点，在经典上都对应于本讲所描述的一个 $U(1)$ 理论。经典真空空间上的度量，由「平凡联络与复化李代数中协变常标量所构成的空间」上的自然度量诱导而来；而后者无非是复化李代数上的 ad 不变度量。限制到上面的对角矩阵，我们看到诱导度量为 $da\wedge d\bar a$。由于 $u=a^2$，这意味着 $u$ 平面上的度量是

\[
  ds^2=\frac{du\,d\bar u}{|u|^2}
\]

的倍数，

% [Tutor Note: 这里的分母疑为笔误。按本页给出的 $da\wedge d\bar a$ 与 $u=a^2$ 换算：$du=2a\,da$，故 $da\,d\bar a=\frac{du\,d\bar u}{4|u|}$，$u$ 平面上的度量应是 $\frac{du\,d\bar u}{|u|}$ 的倍数，分母只有一个 $|u|$。量纲也支持这一点：$u$ 的量纲为 $2$、$ds^2$ 的量纲为 $2$，而 $|du|^2/|u|^2$ 无量纲。几何上二者差别不小：$\frac{|du|^2}{|u|}$ 在原点给出锥（令 $u=re^{i\theta}$，得 $\frac{dr^2}{r}+r\,d\theta^2$，再置 $\rho=2\sqrt r$ 化为 $d\rho^2+\frac{\rho^2}{4}d\theta^2$，即亏欠角为 $\pi$ 的锥），恰是 Weyl 群 $\mathbf{Z}_2$（$a\to-a$）商所应有的锥奇点；而 $\frac{|du|^2}{|u|^2}$ 在原点附近形如柱面端，并不给出锥。下一句「在原点有奇点」所指的正是这个锥奇点。原文只说 $ds^2$ 是「一个倍数」，若按 $\frac{|du|^2}{|u|}$ 理解，还须自行补上那个常数倍数（它含 $\mathrm{Im}\,\tau$ 与 $1/4\pi$）。]

% ===== 原文 p.245 =====

这个度量在原点处有奇点。该奇点反映了这样一个事实：在原点处还有别的无质量场，因此在这一点上低能有效理论并不是纯 $\mathcal{N}=2$ 超对称 $U(1)$ 规范理论；若试图在 $u=0$ 附近把它描述为纯 $U(1)$ 理论，就会得到奇点。

现在，在经典理论中，我们可以描述「以 $0$ 为挖孔点（punctured）的 $u$ 平面」上的环面族（family of tori），以及该族全空间上的全纯二形式 $\eta$。函数 $\tau$ 为常数，因此环面族的单值性包含在 $\{\pm1\}\subset SL(2,\mathbf{Z})$ 之中。由于 $u$ 平面去掉原点后的度量由 $\mathrm{Im}\,\tau\,da\wedge d\bar a$ 给出，我们看到：参数 $a=\sqrt u$ 就是 $U(1)$ 理论中同名参数（即底空间上一对互为对偶的局部全纯坐标中的第一个），而对偶参数为 $a_D=\tau a$。公式 $u=a^2$ 或 $a=\sqrt u$ 意味着：$u$ 平面中绕无穷远（等价地，绕零点）的单值性把 $a$ 送到 $-a$，把 $a_D$ 送到 $-a_D$。这意味着：$u$ 平面（去掉原点）上的环面族，可由 $a$ 平面（去掉原点）上的乘积族 $(\mathbf{C}-\{0\})\times E$ 商去对合 $(a,z)\mapsto(-a,-z)$ 而得到。（这里 $\tau$ 是出现在经典 $SU(2)$ 拉格朗日量中的复参数，$E$ 是椭圆曲线 $\mathbf{C}/L(1,\tau)$，即复平面商去由 $1$ 与 $\tau$ 生成的格；$z\mapsto-z$ 是椭圆曲线上的「乘以 $-1$」。）全纯二形式 $\eta$ 是 $da\wedge dz$ 的像，其中 $dz$ 是 $E$ 上的整体全纯一形式。

% [Tutor Note: 此处有两个极易混淆的等式：$a_D=\tau a$ 与 $u=a^2$。前者是经典理论中 $\tau$ 为常数时的对偶坐标关系（不是模变换！），后者是 Weyl 不变量的定义。切勿因为 $a_D=\tau a$ 就把 $a_D$ 当成「按 $SL(2,\mathbf{Z})$ 变换的第二个坐标」：经典理论中 $\tau$ 固定，绕无穷远的单值性只有 $a\to-a$、$a_D\to-a_D$，即 $\{\pm1\}$。真正的 $SL(2,\mathbf{Z})$ 单值性只在量子理论中出现（下一讲），那时 $\tau(a)$ 不再是常数、$a_D=\partial\mathcal{F}/\partial a$，绕奇点的单值变换才成为非平凡的 $SL(2,\mathbf{Z})$ 元素。另注意对合 $(a,z)\mapsto(-a,-z)$ 同时翻转两个因子，若只翻转 $a$ 就会得到错误的商空间（拓扑上不再是原来的椭圆纤维化）。]

以上描述属于经典理论。在下一讲中我们将讨论该理论的量子版本。结果表明，过渡到量子理论时情况变化很大，但这些变化在某些方面并不像人们想象的那么剧烈。例如，量子真空的空间仍将是同一个 $u$ 平面，不过奇点会出现在不同的位置，而且它们将对应于不同类型的对象变为无质量。相应地，奇点处的单值性也会不同。

\subsection{参考文献与交叉引用}
% 原文 References and cross references

关于本讲以及下两讲的基本文献是

N. Seiberg and E. Witten, \emph{Electric--Magnetic duality, monopole condensation,
and confinement in $\mathcal{N}=2$ supersymmetric Yang--Mills theory},
Nucl. Phys. B \textbf{426} (1994), pp. 19--52.

关于 BPS 谱的跳变，二维情形可参看

S. Cecotti and C. Vafa, \emph{On Classification Of $\mathcal{N}=2$ Supersymmetric
Theories}, Commun. Math. Phys. \textbf{158} (1993) 569--644.

关于 $\mathcal{N}=2$ 多重态的详细讨论可见 [I-Homework] 的 Problem ASH5。拉格朗日量 (17.1) 在 [I-Supersolutions, §10.3] 中有详细展开。四维阿贝尔对偶的入门材料包含在第 8 讲第 5 节中。BPS 态在第 12 讲中有详细讨论。

% [Tutor Note: 文献代号 [I-Homework]、[I-Supersolutions] 指本丛书第 I 卷所附超对称习题集与习题解答（参见本书附录 B）。「拉格朗日量 (17.1)」应指第 17 讲靠前处（第 1 部分）已编号的 $\mathcal{N}=2$ $U(1)$ 矢量多重态拉格朗日量，其形式与本文件 p.241 上所印（该页未印编号）的 $\mathcal{L}=\frac{1}{4\pi}\int d^4x\,d^2\theta\,i\tau W_\alpha W^\alpha+\text{c.c.}+\int d^4x\,d^4\theta\,\mathrm{Im}\,\tau\,\bar\Phi\Phi$ 一式相同，读者不必在 p.240–245 上寻找编号为 (17.1) 的公式。]
